\chapter{Présentations des outils et les captures d'écran}
\thispagestyle{MyStyle}
\begin{onehalfspace}

\section{Introduction}

Ce chapitre présente les outils retenus pour réaliser la plateforme de dématérialisation de la gestion des candidatures, ainsi que l'application telle qu'elle a été développée dans le cadre du stage. Le chapitre précédent a fixé la conception UML --- cas d'utilisation, séquences et classes. Il s'agit désormais de relier cette conception à l'environnement technique et aux écrans effectivement mis en œuvre.

Nous commençons par décrire les principaux outils et technologies, chacun dans le rôle qu'il joue dans ce projet. Une vue d'ensemble de l'architecture situe ensuite le navigateur, la passerelle, les microservices et les bases de données PostgreSQL. Enfin, des captures d'écran illustrent le parcours métier : connexion, tableau de bord, candidats, concours, lieux, répartition, convocations et, pour l'administrateur, la gestion des comptes gestionnaires.

\section{Présentation des outils}

Cette section décrit l'environnement logiciel utilisé pour réaliser la plateforme. Chaque outil est présenté dans le rôle qu'il joue dans ce projet. PlantUML, déjà présenté au chapitre~3 (figure~\ref{fig:logo-plantuml}), a servi à produire les diagrammes UML ; il n'est pas repris ici.

\subsection{Environnement de développement}

\begin{figure}[H]
  \centering
  
\includegraphics[width=0.28\textwidth]{windows-logo.png}
  \caption{Logo de Windows}
  \label{fig:logo-windows}
\end{figure}

Le développement a été conduit sous \textbf{Windows}. Ce système d'exploitation a accueilli l'éditeur, le JDK, Maven, PostgreSQL et l'environnement Node.js nécessaire au frontend.

\begin{figure}[H]
  \centering
  
\includegraphics[width=0.22\textwidth]{vscode-logo.jpg}
  \caption{Logo de Visual Studio Code}
  \label{fig:logo-vscode}
\end{figure}

\textbf{Visual Studio Code} a été utilisé comme environnement de développement. Il a permis d'éditer le frontend React (TypeScript), les sources Java des microservices, les fichiers Maven et la documentation du projet, dans un même éditeur.

\begin{figure}[H]
  \centering
  
\includegraphics[width=0.32\textwidth]{java-logo.png}
  \caption{Logo de Java}
  \label{fig:logo-java}
\end{figure}

Le langage \textbf{Java 17} constitue le socle du backend. Les six microservices et la passerelle API sont écrits en Java 17 et s'exécutent sur la machine virtuelle Java.

\begin{figure}[H]
  \centering
  
\includegraphics[width=0.42\textwidth]{maven-logo.png}
  \caption{Logo de Maven}
  \label{fig:logo-maven}
\end{figure}

\textbf{Maven} assure la construction du backend, organisé en projet multi-modules : la passerelle et les six services sont déclarés dans un POM parent, qui fixe notamment Spring Boot 3.4.4 et Java 17.

\subsection{Technologies côté serveur}

\begin{figure}[H]
  \centering
  
\includegraphics[width=0.22\textwidth]{spring-boot-logo.png}
  \caption{Logo de Spring Boot}
  \label{fig:logo-spring-boot}
\end{figure}

\textbf{Spring Boot 3.4.4} est le cadre d'application de chaque microservice. Il fournit le serveur web, l'injection de dépendances, l'accès aux données (JPA) et la configuration, de sorte que chaque service reste une application autonome.

\begin{figure}[H]
  \centering
  
\includegraphics[width=0.28\textwidth]{spring-cloud-logo.jpg}
  \caption{Logo de Spring Cloud Gateway}
  \label{fig:logo-spring-cloud}
\end{figure}

\textbf{Spring Cloud Gateway} (Spring Cloud 2024.0.1) constitue le point d'entrée unique des appels HTTP. Le navigateur n'interroge pas les microservices directement : les requêtes transitent par la passerelle, qui les achemine vers le service concerné.

\begin{figure}[H]
  \centering
  
\includegraphics[width=0.42\textwidth]{spring-security-logo.jpg}
  \caption{Logo de Spring Security}
  \label{fig:logo-spring-security}
\end{figure}

\textbf{Spring Security}, associé à des jetons \textbf{JWT}, protège les API. Après une connexion réussie, un jeton d'accès est délivré ; il n'y a pas de jeton de rafraîchissement. Seul un utilisateur authentifié accède aux fonctions autorisées selon son rôle (gestionnaire ou administrateur).

\begin{figure}[H]
  \centering
  
\includegraphics[width=0.28\textwidth]{postgresql-logo.png}
  \caption{Logo de PostgreSQL}
  \label{fig:logo-postgresql}
\end{figure}

\textbf{PostgreSQL} est le système de gestion de bases de données relationnelles. Chaque microservice dispose de \textbf{sa propre base}, ce qui isole les données d'authentification, des candidats, des concours, des lieux, de la répartition et des envois de convocations.

\begin{figure}[H]
  \centering
  
\includegraphics[width=0.28\textwidth]{flyway-logo.png}
  \caption{Logo de Flyway}
  \label{fig:logo-flyway}
\end{figure}

\textbf{Flyway} gère les migrations de schéma : chaque service applique ses propres scripts à sa base PostgreSQL, de manière reproductible au démarrage.

\begin{figure}[H]
  \centering
  
\includegraphics[width=0.48\textwidth]{apache-poi-logo.png}
  \caption{Logo d'Apache POI}
  \label{fig:logo-apache-poi}
\end{figure}

\textbf{Apache POI} (version 5.3.0) est utilisé dans le service candidats pour lire les fichiers Excel (\texttt{.xlsx}) lors de l'import de la liste. Les candidats ne sont pas créés un par un par l'interface : l'entrée dans le système se fait par cet import.

\begin{figure}[H]
  \centering
  
\includegraphics[width=0.22\textwidth]{openpdf-logo.png}
  \caption{Logo d'OpenPDF}
  \label{fig:logo-openpdf}
\end{figure}

\textbf{OpenPDF} (version 1.3.30) génère, côté serveur, le PDF unitaire de convocation, à la demande, pour un candidat dont l'affectation est complète (centre, établissement, salle et place).

\begin{figure}[H]
  \centering
  
\includegraphics[width=0.42\textwidth]{spring-mail-logo.png}
  \caption{Logo de Spring Mail}
  \label{fig:logo-spring-mail}
\end{figure}

\textbf{Spring Mail} envoie les convocations par courrier électronique, via le SMTP de \textbf{Gmail}. Sans identifiants de messagerie configurés, le service démarre mais l'envoi est refusé.

\subsection{Technologies côté client}

\begin{figure}[H]
  \centering
  
\includegraphics[width=0.28\textwidth]{react-logo.png}
  \caption{Logo de React}
  \label{fig:logo-react}
\end{figure}

\textbf{React 19}, avec \textbf{TypeScript}, constitue l'interface web : une application monopage qui présente les écrans de connexion, de tableau de bord, de candidats, de concours, de lieux, de répartition, de convocations et de gestion des comptes. La navigation entre ces pages est assurée par React Router.

\begin{figure}[H]
  \centering
  
\includegraphics[width=0.22\textwidth]{typescript-logo.png}
  \caption{Logo de TypeScript}
  \label{fig:logo-typescript}
\end{figure}

\textbf{TypeScript} apporte un typage statique au code de l'interface. Il réduit les erreurs de contrat entre les pages React et les réponses JSON des API, notamment pour les concours, les candidats et les convocations.

\begin{figure}[H]
  \centering
  
\includegraphics[width=0.42\textwidth]{vite-logo.png}
  \caption{Logo de Vite}
  \label{fig:logo-vite}
\end{figure}

\textbf{Vite 6} sert d'outil de construction et de serveur de développement du frontend. Il compile le projet React et TypeScript et fournit le bundle exécuté dans le navigateur.

\begin{figure}[H]
  \centering
  
\includegraphics[width=0.42\textwidth]{axios-logo.png}
  \caption{Logo d'axios}
  \label{fig:logo-axios}
\end{figure}

\textbf{axios} est la bibliothèque HTTP du frontend. Toutes les requêtes de l'interface --- authentification, import Excel, concours, lieux, répartition, convocations --- sont envoyées à la passerelle par son intermédiaire.

\begin{figure}[H]
  \centering
  
\includegraphics[width=0.42\textwidth]{jspdf-logo.png}
  \caption{Logo de jsPDF}
  \label{fig:logo-jspdf}
\end{figure}

\textbf{jsPDF} permet d'exporter, depuis le navigateur, des listes au format PDF (résultats de répartition et convocations). Ces exports de listes sont distincts du PDF unitaire de convocation produit par OpenPDF côté serveur.

\begin{figure}[H]
  \centering
  
\includegraphics[width=0.22\textwidth]{docx-logo.png}
  \caption{Logo de la bibliothèque docx}
  \label{fig:logo-docx}
\end{figure}

La bibliothèque \textbf{docx} permet le même type d'export au format Word, toujours côté client, sur les pages de répartition et de convocations.

\section{Architecture technique}

La plateforme est organisée en une application web, exécutée dans le navigateur, et un backend en microservices Spring Boot, exposés derrière une passerelle unique. Cette organisation isole chaque domaine métier --- authentification, candidats, concours, lieux, répartition, convocations --- tout en n'offrant qu'un seul point d'entrée aux écrans. Elle met en œuvre, côté technique, la conception UML du chapitre~3.

\subsection{Vue d'ensemble}

L'utilisateur travaille dans le navigateur. L'interface React envoie ses requêtes HTTP à la \textbf{passerelle API} (Spring Cloud Gateway). Celle-ci n'héberge pas de données et ne vérifie pas le jeton : elle achemine chaque appel vers le microservice correspondant, d'après le chemin de l'URL, et \textbf{relaye} l'en-tête d'autorisation. Le service cible traite la demande, interroge sa base PostgreSQL le cas échéant, et renvoie la réponse par le même chemin.

Certains traitements croisent plusieurs domaines. La répartition consulte les concours, les salles et les candidats, puis enregistre les affectations. Les convocations sont \textbf{assemblées à la volée} à partir des candidats, des concours et des lieux. Ces appels entre services transmettent le jeton JWT déjà obtenu à la connexion.

\subsection{Microservices et persistance}

Le backend comprend \textbf{six microservices}, plus la passerelle :

\begin{table}[H]
\centering
\caption{Composants du backend}
\label{tab:composants-backend}
\begin{spacing}{1}
\resizebox{\textwidth}{!}{%
\begin{tabular}{|p{5.2cm}|p{10.2cm}|}
\hline
\rowcolor[HTML]{C9DAF8}
\textbf{Composant} & \textbf{Rôle} \\
\hline
Passerelle API &
Point d'entrée unique et routage des requêtes \\
\hline
Service d'authentification &
Connexion, délivrance du jeton JWT, comptes utilisateurs \\
\hline
Service candidats &
Import Excel, consultation, mise à jour, suppression et affectations \\
\hline
Service concours &
Concours et centres d'examen associés \\
\hline
Service lieux &
Centres, établissements et salles \\
\hline
Service répartition &
Répartition automatique et historique des exécutions \\
\hline
Service convocations &
Assemblage des convocations, PDF unitaire, envoi par e-mail et historique des envois \\
\hline
\end{tabular}%
}
\end{spacing}
\end{table}

Chaque microservice dispose de \textbf{sa propre base PostgreSQL}. Les liens entre domaines (numéro d'inscription, numéro de concours, identifiants de centre, d'établissement et de salle) sont des \textbf{références logiques}, contrôlées par appels HTTP, et non des clés étrangères partagées entre bases. Une convocation n'est pas une table : seul l'historique des envois est persisté.

\subsection{Authentification et rôles}

L'authentification s'effectue par \texttt{POST /auth/login}, via la passerelle. En cas de succès, un \textbf{jeton d'accès JWT} est délivré ; il n'existe pas de jeton de rafraîchissement. L'interface le conserve dans le \texttt{sessionStorage} du navigateur et l'ajoute ensuite aux requêtes. \textbf{Chaque service ressource} vérifie le jeton et le rôle : le gestionnaire dispose des droits d'écriture métier ; l'administrateur consulte le métier en lecture seule et gère les comptes gestionnaires.

\section{Interfaces de l'application}

Cette section illustre l'application réalisée, à l'aide de captures correspondant au parcours du \textbf{gestionnaire} : authentification, tableau de bord, puis les modules métier. Les boutons d'écriture (import, création, lancement de la répartition) n'apparaissent que pour ce rôle.

\subsection{Authentification}

\begin{figure}[H]
  \centering
  \includegraphics[width=\textwidth]{screenshot-connexion.png}
  \caption{Page de connexion}
  \label{fig:screenshot-connexion}
\end{figure}

La figure présente l'écran d'accès à l'application. Le bandeau affiche le logo du \textbf{Ministère de l'Économie et des Finances}, le titre «~Connexion~» et le sous-titre «~Gestion des candidatures --- authentification~». Le gestionnaire saisit son nom d'utilisateur et son mot de passe, puis valide par «~Se connecter~». Aucun autre écran n'est accessible tant que l'authentification n'a pas abouti.

\subsection{Tableau de bord}

\begin{figure}[H]
  \centering
  \includegraphics[width=\textwidth,height=0.88\textheight,keepaspectratio]{screenshot-dashboard.png}
  \caption{Tableau de bord}
  \label{fig:screenshot-dashboard}
\end{figure}

La figure présente l'écran d'accueil après connexion. L'en-tête reprend le logo du ministère, les onglets (Tableau de bord, Candidats, Concours, Lieux, Répartition, Convocations) et le rôle \textbf{Gestionnaire}, avec le bouton «~Déconnexion~». Le gestionnaire dispose d'un bouton «~Actualiser~». Six indicateurs résument l'état : 104 candidats inscrits, 104 affectés, 6 concours, 4 centres, une capacité de 720 places et 104 convocations (dont 100 envoyées et 4 échecs). La barre «~Taux d'affectation global~» indique que les 104 candidats ont un centre, un établissement, une salle et une place. La colonne de gauche détaille les effectifs par concours ; celle de droite rappelle la dernière répartition, au statut «~Terminée~», déclenchée par le gestionnaire.

\subsection{Gestion des candidats}

\begin{figure}[H]
  \centering
  \includegraphics[width=\textwidth,height=0.88\textheight,keepaspectratio]{screenshot-candidats.png}
  \caption{Liste des candidats}
  \label{fig:screenshot-candidats}
\end{figure}

La figure présente l'écran «~Liste des candidats~». Le gestionnaire peut \textbf{importer un fichier Excel} ou \textbf{réinitialiser la liste} ; il n'existe pas de création unitaire de candidat. Des filtres permettent de rechercher par nom, prénom, CIN ou numéro d'inscription, et de restreindre par ville ou par concours. Chaque fiche affiche l'identité, les coordonnées, la ville, le concours, puis l'affectation (centre, établissement, salle et place). Les actions «~Modifier~» et «~Supprimer~» figurent en bas de fiche.

\subsection{Gestion des concours}

\begin{figure}[H]
  \centering
  \includegraphics[width=\textwidth,height=0.88\textheight,keepaspectratio]{screenshot-concours-liste.png}
  \caption{Liste des concours}
  \label{fig:screenshot-concours-liste}
\end{figure}

La figure présente l'écran «~Concours~». Le gestionnaire actualise la liste ou ouvre la création par «~+ Nouveau concours~». Trois indicateurs donnent le nombre de concours enregistrés, les examens à venir et les liens vers les centres. Chaque carte porte le nom, le numéro, la date et l'heure d'examen, les \textbf{centres assignés} (Tanger, Marrakech, Casablanca, Rabat sur cette vue) et les boutons «~Modifier~» et «~Supprimer~».

\begin{figure}[H]
  \centering
  \includegraphics[width=\textwidth]{screenshot-concours-creation.png}
  \caption{Création d'un concours}
  \label{fig:screenshot-concours-creation}
\end{figure}

La figure présente la fenêtre ouverte par «~+ Nouveau concours~». Le gestionnaire renseigne le nom, le numéro, la date et l'heure d'examen, puis sélectionne \textbf{au moins un centre} parmi ceux déjà créés dans les lieux. Il peut ajouter d'autres centres. «~Enregistrer~» envoie la création ; «~Annuler~» ferme la fenêtre sans enregistrement.

\subsection{Gestion des lieux}

\begin{figure}[H]
  \centering
  \includegraphics[width=\textwidth,height=0.88\textheight,keepaspectratio]{screenshot-lieux.png}
  \caption{Gestion des lieux d'examen}
  \label{fig:screenshot-lieux}
\end{figure}

La figure présente l'écran «~Lieux~». À gauche, la liste des \textbf{centres} (Casablanca, Marrakech, Rabat, Tanger) ; Casablanca est sélectionné. Le détail montre la hiérarchie \textbf{centre $\rightarrow$ établissement $\rightarrow$ salles} : l'établissement Ibn Tofail et des salles (A, B, C) avec leur capacité et le concours associé. Le gestionnaire dispose des boutons «~+ Nouveau centre~», «~+ Établissement~», «~+ Salle~», ainsi que «~Modifier~» et «~Supprimer~» à chaque niveau.

\subsection{Répartition automatique}

\begin{figure}[H]
  \centering
  \includegraphics[width=\textwidth,height=0.88\textheight,keepaspectratio]{screenshot-repartition-lancement.png}
  \caption{Lancement de la répartition automatique}
  \label{fig:screenshot-repartition-lancement}
\end{figure}

La figure présente l'écran «~Répartition~». Le gestionnaire peut actualiser, \textbf{réinitialiser} les sièges (sans vider la liste des candidats) ou \textbf{lancer la répartition}. Après le clic sur «~Lancer la répartition~», une confirmation demande s'il faut exécuter une nouvelle répartition automatique de tous les candidats. En arrière-plan, l'historique des exécutions affiche le statut («~Terminée~» ou «~Terminée avec alertes~»), le taux d'affectation et les boutons «~Voir la synthèse~» et «~Supprimer~».

\begin{figure}[H]
  \centering
  \includegraphics[width=\textwidth,height=0.88\textheight,keepaspectratio]{screenshot-repartition-synthese.png}
  \caption{Synthèse d'une exécution de répartition}
  \label{fig:screenshot-repartition-synthese}
\end{figure}

La figure présente la fenêtre «~Synthèse · exécution \#17~», ouverte depuis «~Voir la synthèse~». L'exécution, déclenchée par le gestionnaire, est au statut «~Terminée~» : 104 candidats, 104 affectés, aucune alerte, taux d'affectation de 100~\%. Le tableau liste, pour chaque candidat, la ville, le concours, le centre, l'établissement, la salle et la place. Des filtres par centre et par concours, ainsi que les boutons \textbf{Word} et \textbf{PDF}, permettent d'exporter la liste \textbf{depuis le navigateur} (distinct du PDF unitaire de convocation produit côté serveur).

\begin{figure}[H]
  \centering
  \includegraphics[width=\textwidth,height=0.88\textheight,keepaspectratio]{screenshot-repartition-export.png}
  \caption{Export PDF d'une synthèse filtrée}
  \label{fig:screenshot-repartition-export}
\end{figure}

La figure présente à nouveau la synthèse de l'exécution n\textsuperscript{o}~17, après filtrage. Le gestionnaire a choisi le centre \textbf{Casablanca} ; le compteur indique \textbf{24 affectations sur 104}. Le tableau ne retient que les lignes de ce centre (établissement, salle, place). Deux groupes d'export apparaissent : \textbf{Exporter la sélection} (Word ou PDF, limité au filtre) et \textbf{Toutes les affectations} (liste complète). Le fichier téléchargé, \texttt{repartition-17-filtre.pdf}, correspond à la sélection. Cet export est produit \textbf{dans le navigateur} ; il ne se confond pas avec le PDF unitaire de convocation généré côté serveur par OpenPDF.

\subsection{Gestion des convocations}

\begin{figure}[H]
  \centering
  \includegraphics[width=\textwidth,height=0.88\textheight,keepaspectratio]{screenshot-convocations.png}
  \caption{Page des convocations}
  \label{fig:screenshot-convocations}
\end{figure}

La figure présente l'écran «~Convocations~». L'en-tête identifie le \textbf{gestionnaire}. Celui-ci peut actualiser la page ou cliquer sur \textbf{Envoyer toutes les convocations}. Quatre indicateurs résument l'état : 104 convocations prêtes (candidats affectés), 104 avec e-mail, 0 sans e-mail, 100 envois réussis (4 échecs). Le texte précise qu'une convocation est assemblée pour chaque candidat dont l'affectation est complète ; il n'y a pas de bouton «~Générer~». Des filtres par centre et par concours, une barre \textbf{Exporter la liste} (Word et PDF) et une grille de cartes (identité, concours, centre, établissement, salle, place, date, e-mail, \textbf{Voir le PDF}) composent le reste de l'écran.

\begin{figure}[H]
  \centering
  \includegraphics[width=\textwidth,height=0.88\textheight,keepaspectratio]{screenshot-convocation-pdf.png}
  \caption{PDF unitaire de convocation}
  \label{fig:screenshot-convocation-pdf}
\end{figure}

La figure présente le document obtenu par \textbf{Voir le PDF} pour un candidat. L'en-tête indique «~Royaume du Maroc --- Organisation des concours~», le titre \textbf{CONVOCATION À L'EXAMEN} et le nom du concours. Le corps convoque le candidat aux date, lieu et place indiqués, puis un tableau reprend le nom, le numéro d'inscription, le concours, la date et l'heure, le centre, l'établissement, la salle et la place. Un encadré rappelle de se présenter en avance, muni de la convocation et d'une pièce d'identité. Ce fichier est produit \textbf{côté serveur} par OpenPDF, à la demande, pour un candidat déjà affecté.

\begin{figure}[H]
  \centering
  \includegraphics[width=\textwidth,height=0.88\textheight,keepaspectratio]{screenshot-convocations-export.png}
  \caption{Export PDF de la liste des convocations}
  \label{fig:screenshot-convocations-export}
\end{figure}

La figure présente le même écran, après un clic sur \textbf{PDF} dans «~Exporter la liste~». Le navigateur télécharge le fichier \texttt{convocations-complet.pdf}. Contrairement à la figure précédente, il s'agit d'une \textbf{liste} de convocations, générée \textbf{dans le navigateur}, et non du document unitaire OpenPDF.

\begin{figure}[H]
  \centering
  \includegraphics[width=\textwidth,height=0.88\textheight,keepaspectratio]{screenshot-convocations-liste-pdf.png}
  \caption{Liste des convocations exportée}
  \label{fig:screenshot-convocations-liste-pdf}
\end{figure}

La figure présente le contenu de cet export. Le titre est \textbf{Liste des convocations} ; l'en-tête indique la date d'export et le périmètre (ici : toutes les convocations, 104). Un tableau à colonnes (candidat, numéro d'inscription, concours, centre, établissement, salle, date et heure, place, e-mail) occupe le corps du document. Ce PDF de liste est distinct de la convocation nominative de la figure~\ref{fig:screenshot-convocation-pdf}.

\begin{figure}[H]
  \centering
  \includegraphics[width=\textwidth]{screenshot-convocations-envoi.png}
  \caption{Résultat de l'envoi groupé des convocations}
  \label{fig:screenshot-convocations-envoi}
\end{figure}

La figure présente la fenêtre ouverte après \textbf{Envoyer toutes les convocations}. Le titre est «~Résultat de l'envoi~», le sous-titre «~Récapitulatif de l'envoi groupé par e-mail~». Trois totaux (Total, Envoyés, Échecs) résument l'opération ; sur cette capture, 104 envois aboutis, aucun échec. La liste détaille chaque destinataire (nom, adresse, statut \textbf{Envoyé}). «~Fermer~» revient à la page des convocations.

\subsection{Vue administrateur (lecture seule sur le métier)}

\begin{figure}[H]
  \centering
  \includegraphics[width=\textwidth,height=0.88\textheight,keepaspectratio]{screenshot-admin-candidats.png}
  \caption{Liste des candidats en lecture seule}
  \label{fig:screenshot-admin-candidats}
\end{figure}

La figure présente le même écran «~Liste des candidats~», ouvert avec le compte \textbf{administrateur}. L'en-tête affiche le logo du ministère, l'onglet \textbf{Gestionnaires} (absent chez le gestionnaire) et la pastille \textbf{Administrateur · lecture seule}. Les filtres (recherche, ville, concours) et les fiches restent visibles, y compris l'affectation. En revanche, les boutons \textbf{Importer fichier Excel}, \textbf{Réinitialiser la liste}, \textbf{Modifier} et \textbf{Supprimer} n'apparaissent pas : l'administrateur consulte, il ne modifie pas les candidatures.

\begin{figure}[H]
  \centering
  \includegraphics[width=\textwidth,height=0.88\textheight,keepaspectratio]{screenshot-admin-concours.png}
  \caption{Liste des concours en lecture seule}
  \label{fig:screenshot-admin-concours}
\end{figure}

La figure présente l'écran «~Concours~» pour le même rôle. Le bouton «~Actualiser~» est présent ; «~+ Nouveau concours~» n'apparaît pas, pas plus que «~Modifier~» ou «~Supprimer~» sur les cartes. L'administrateur voit le nom, le numéro, la date et les centres assignés, sans pouvoir créer ni altérer un concours.

\subsection{Gestion des comptes gestionnaires}

\begin{figure}[H]
  \centering
  \includegraphics[width=\textwidth,height=0.88\textheight,keepaspectratio]{screenshot-gestionnaires-liste.png}
  \caption{Liste des gestionnaires}
  \label{fig:screenshot-gestionnaires-liste}
\end{figure}

La figure présente l'écran réservé à l'administrateur : «~Liste des gestionnaires~». Le tableau indique le nom d'utilisateur, le statut (ici \textbf{Actif}), la date de création et les actions \textbf{Modifier} et \textbf{Supprimer}. Le bouton \textbf{Nouveau gestionnaire} ouvre la création d'un compte. Seul l'administrateur accède à cette page ; le gestionnaire ne la voit pas dans le menu.

\begin{figure}[H]
  \centering
  \includegraphics[width=\textwidth]{screenshot-gestionnaires-creation.png}
  \caption{Création d'un compte gestionnaire}
  \label{fig:screenshot-gestionnaires-creation}
\end{figure}

La figure présente la fenêtre ouverte par «~Nouveau gestionnaire~». L'administrateur saisit un nom d'utilisateur et un mot de passe. Une mention indique que le rôle est fixé à \texttt{GESTIONNAIRE} et que le mot de passe doit compter au moins huit caractères. «~Enregistrer~» crée le compte ; «~Annuler~» ferme la fenêtre sans enregistrement.

\begin{figure}[H]
  \centering
  \includegraphics[width=\textwidth]{screenshot-gestionnaires-modification.png}
  \caption{Modification d'un compte gestionnaire}
  \label{fig:screenshot-gestionnaires-modification}
\end{figure}

La figure présente la fenêtre ouverte par «~Modifier~». Le nom d'utilisateur peut être mis à jour. Le champ «~Nouveau mot de passe~» peut rester vide pour conserver le mot de passe actuel. La case \textbf{Compte actif (peut se connecter)} permet d'activer ou de désactiver l'accès. Le rôle demeure \texttt{GESTIONNAIRE}. «~Enregistrer~» valide les changements.

\section{Conclusion}

Ce chapitre a présenté l'environnement technique de la plateforme de dématérialisation de la gestion des candidatures, ainsi que l'application telle qu'elle a été développée dans le cadre du stage. Les outils retenus --- Java 17, Maven, Spring Boot, Spring Cloud Gateway, Spring Security et JWT, PostgreSQL, Flyway, Apache POI, OpenPDF, Spring Mail, React, TypeScript, Vite, axios, ainsi que jsPDF et docx --- ont été situés dans le rôle qu'ils jouent dans ce projet. L'architecture a ensuite été exposée : le navigateur s'adresse à une passerelle unique, qui achemine les requêtes vers six microservices, chacun disposant de sa propre base PostgreSQL ; le jeton d'accès JWT est relayé et vérifié par chaque service ressource.

Les captures d'écran ont illustré le parcours réel. Le gestionnaire s'authentifie, consulte le tableau de bord, importe les candidats par fichier Excel, organise les concours et les lieux, lance la répartition automatique, consulte la synthèse des affectations, exporte les listes, visualise le PDF unitaire de convocation et procède à l'envoi groupé par e-mail. L'administrateur dispose des mêmes vues métier en lecture seule et reste seul à gérer les comptes gestionnaires. Deux productions PDF ont été distinguées : le document nominatif, généré côté serveur par OpenPDF, et les listes exportées depuis le navigateur.

Cette présentation relie la conception UML du chapitre précédent à la solution réalisée. Elle clôt la partie technique du rapport et ouvre sur la conclusion générale, les perspectives et les références.

\end{onehalfspace}
